\documentclass[11pt,a4paper]{article}
\usepackage[utf8]{inputenc}
\usepackage[spanish]{babel}
\usepackage{geometry}
\geometry{left=2.5cm, right=2.5cm, top=3cm, bottom=3cm}
\usepackage{tabularx}
\usepackage{booktabs}
\usepackage{xcolor}
\usepackage{colortbl}
\usepackage{hyperref}
\usepackage{titlesec}
\usepackage{enumitem}

% Colores
\definecolor{uteqblue}{RGB}{20, 50, 100}
\definecolor{reqheader}{RGB}{235, 240, 245}

% Estilo de hipervínculos
\hypersetup{
    colorlinks=true,
    linkcolor=uteqblue,
    urlcolor=uteqblue
}

\begin{document}

% Portada
\begin{titlepage}
    \centering
    \vspace*{1cm}
    {\LARGE \textbf{UNIVERSIDAD TÉCNICA ESTATAL DE QUEVEDO}}\\[0.5cm]
    {\large Facultad de Ciencias de la Ingeniería}\\[0.2cm]
    {\large Carrera de Ingeniería de Software}\\[0.8cm]
    
    {\large ASIGNATURA: Aplicaciones Web}\\[0.2cm]
    {\large Quinto Nivel — Período Académico 2026-2027}\\[1cm]
    
    {\Huge \bfseries Especificación de Requisitos de Software (SRS)}\\[0.5cm]
    {\large Estándar ISO/IEC/IEEE 29148:2018}\\[1cm]
    
    {\huge \bfseries \color{uteqblue} Sistema de Simulación de Comportamiento Vial con Inteligencia Artificial (SBVIA)}\\[0.5cm]
    {\Large Versión 1.1.0}\\[1cm]
    
    \begin{table}[h]
        \centering
        \renewcommand{\arraystretch}{1.5}
        \begin{tabularx}{\textwidth}{|l|X|}
            \hline
            \textbf{Repositorio} & \url{https://github.com/gleiston-guerrero/SBVIA-AppWeb} \\ \hline
            \textbf{Punto del historial} & etiqueta v1.1.0 \\ \hline
            \textbf{DOI del software (Zenodo)} & \texttt{10.5281/zenodo.22740480} \\ \hline
            \textbf{Fecha de emisión} & 18 de septiembre de 2026 \\ \hline
        \end{tabularx}
    \end{table}

    \vfill

    {\Large \textbf{AUTORES}}\\[0.2cm]
    {\large Cruz Pérez, Justyn K.}\\[0.1cm]
    {\normalsize ORCID: \texttt{0009-0009-0847-0780}}\\[0.1cm]
    {\normalsize Correo: \texttt{jcruzp@uteq.edu.ec}}\\[0.5cm]

    {\large Umaginga Arévalo, Jefferson M.}\\[0.1cm]
    {\normalsize ORCID: \texttt{0009-0005-9831-0801}}\\[0.1cm]
    {\normalsize Correo: \texttt{jumagingaa@uteq.edu.ec}}\\[0.5cm]

    {\Large \textbf{DOCENTE-DIRECTOR}}\\[0.2cm]
    {\large Ing. Gleiston Cíceron Guerrero Ulloa, Ph.D.}\\[0.5cm]
    
\end{titlepage}

\newpage
\tableofcontents
\newpage

\section*{Control del documento}
\addcontentsline{toc}{section}{Control del documento}

\begin{table}[h]
    \centering
    \renewcommand{\arraystretch}{1.5}
    \begin{tabularx}{\textwidth}{@{}|l|l|X|l|@{}}
        \hline
        \rowcolor{uteqblue}
        \textcolor{white}{\textbf{Versión}} & \textcolor{white}{\textbf{Fecha}} & \textcolor{white}{\textbf{Descripción}} & \textcolor{white}{\textbf{Estado}} \\ \hline
        1.0.0 & 17-09-2026 & Línea base de la entrega final & Aprobada \\ \hline
        1.1.0 & 18-09-2026 & Correcciones finales (contrato REST, Lighthouse, URLs y limpieza) & Pendiente de aprobación \\ \hline
    \end{tabularx}
\end{table}

\newpage

% ==============================
% 1. INTRODUCCIÓN
% ==============================
\section{Introducción}

\subsection{Propósito}
Este documento establece los requisitos verificables de SBVIA v1.1.0 y constituye la línea base contractual para desarrollo, prueba, despliegue y evaluación académica.

\subsection{Alcance}
SBVIA es una aplicación web de formación vial que administra usuarios y escenarios, conserva prácticas de simulación, calcula resultados y presenta retroalimentación. Incluye un backend Java 21 con Spring Boot, un frontend Angular, PostgreSQL 16, Redis 7 y orquestación Docker Compose.

\subsection{Audiencia}
Equipo de desarrollo, docente-director, evaluadores, administradores, instructores y conductores en formación.

\subsection{Definiciones}
\begin{table}[h]
    \centering
    \renewcommand{\arraystretch}{1.3}
    \begin{tabularx}{\textwidth}{@{}|l|X|@{}}
        \hline
        \rowcolor{uteqblue}
        \textcolor{white}{\textbf{Término}} & \textcolor{white}{\textbf{Definición}} \\ \hline
        CRUD-ORM & Operaciones elementales implementadas mediante Spring Data JPA. \\ \hline
        SP & Procedimiento almacenado o función SQL parametrizada. \\ \hline
        JWT & Token firmado para autenticación stateless. \\ \hline
        SUS & System Usability Scale. \\ \hline
        p95 & Percentil 95 de una distribución de latencias. \\ \hline
        VERIFIED & El requisito ha sido implementado y comprobado mediante pruebas automatizadas exhaustivas o auditorías. \\ \hline
        TESTED & El requisito ha sido verificado únicamente a nivel de componente o prueba superficial, o implementado exitosamente. \\ \hline
        IMPLEMENTED & El requisito está desarrollado en el código fuente, pero su cobertura de prueba no valida completamente el criterio. \\ \hline
        Inactividad & Ausencia de autenticación y uso del sistema por un período superior a 6 meses. \\ \hline
        Práctica & Registro individual del desempeño de un conductor en un escenario específico. \\ \hline
        Certificado & Documento inmutable emitido al aprobar un escenario. \\ \hline
        Puntaje & Valor numérico de 0 a 100 asignado por la simulación. \\ \hline
        Escenario & Configuración específica del entorno vial (clima, tráfico, dificultad) a ser simulado. \\ \hline
        Infracción & Violación a una regla de tránsito registrada durante una práctica. \\ \hline
        Comportamiento Vial & Decisión tomada por el usuario evaluada por el motor de IA. \\ \hline
    \end{tabularx}
\end{table}

% ==============================
% 2. DESCRIPCIÓN GENERAL
% ==============================
\section{Descripción general}

\subsection{Perspectiva del producto}
La solución adopta arquitectura cliente-servidor. Angular consume una API REST protegida por JWT; Spring Boot implementa lógica, autorización por roles y persistencia; PostgreSQL conserva los datos y Redis gestiona caché y revocación de tokens.

\subsection{Clases de usuario}
\begin{table}[h]
    \centering
    \renewcommand{\arraystretch}{1.3}
    \begin{tabularx}{\textwidth}{@{}|l|X|@{}}
        \hline
        \rowcolor{uteqblue}
        \textcolor{white}{\textbf{Rol}} & \textcolor{white}{\textbf{Capacidades principales}} \\ \hline
        Conductor & Registro, autenticación, consulta de escenarios, prácticas, métricas y perfil. \\ \hline
        Administrador & Gestión de escenarios, usuarios, roles y consulta global de prácticas. \\ \hline
        Instructor & Consulta de resultados agregados de estudiantes a su cargo. \\ \hline
    \end{tabularx}
\end{table}

\subsection{Restricciones}
La versión estable utiliza Java 21, Spring Boot, Angular 17 o superior, PostgreSQL 16, Flyway, Redis 7 y Docker Compose. No se admiten secretos productivos versionados. El despliegue público debe usar HTTPS.

\subsection{Suposiciones y dependencias}
El navegador admite cookies SameSite y JavaScript moderno. El entorno dispone de Docker y conectividad entre servicios. Los servicios externos de publicación, DOI y dominio son responsabilidad organizativa y no se simulan dentro del software.

\newpage
% ==============================
% 3. REQUISITOS FUNCIONALES
% ==============================
\section{Requisitos funcionales}

\newenvironment{requisito}[1]{
    \vspace{0.5cm}
    \noindent\tabularx{\textwidth}{@{}|X|@{}}
    \hline
    \rowcolor{reqheader} \textbf{#1} \\
    \hline
}{
    \\ \hline
    \endtabularx
}

\begin{requisito}{RF-01 - Autenticación stateless con JWT}
    \textbf{Tipo:} FUNC \quad \textbf{Prioridad:} Must \quad \textbf{Estado:} VERIFIED \\[0.2cm]
    \textbf{Módulo:} Backend \quad \textbf{Acceso a datos:} CRUD-ORM \\[0.2cm]
    \textbf{Descripción:} El sistema deberá autenticar a los usuarios mediante credenciales válidas y emitir un JWT (JSON Web Token) stateless. \\[0.2cm]
    \textbf{Criterio de aceptación:} Dado un usuario registrado, cuando presenta credenciales válidas, entonces recibe un JWT vigente y cookies seguras; credenciales inválidas producen 401. \\[0.2cm]
    \textbf{Verificación:} Prueba de Integración (AuthServiceTest.java) \quad \textbf{SQL:} No aplica
\end{requisito}

\begin{requisito}{RF-02 - Gestión de Escenarios viales (CRUD)}
    \textbf{Tipo:} FUNC \quad \textbf{Prioridad:} Must \quad \textbf{Estado:} VERIFIED \\[0.2cm]
    \textbf{Módulo:} Backend \quad \textbf{Acceso a datos:} CRUD-ORM \\[0.2cm]
    \textbf{Descripción:} El sistema deberá permitir a los administradores crear, leer, actualizar y eliminar (desactivar) escenarios viales, restringiendo la escritura a los conductores. \\[0.2cm]
    \textbf{Criterio de aceptación:} Un administrador puede crear, consultar, actualizar y desactivar escenarios; un conductor solo puede consultarlos y recibe 403 al intentar modificarlos. \\[0.2cm]
    \textbf{Verificación:} Prueba de Integración (EscenarioServiceTest.java) \quad \textbf{SQL:} No aplica
\end{requisito}

\begin{requisito}{RF-03 - Consulta multi-tabla de resultados de simulación}
    \textbf{Tipo:} FUNC \quad \textbf{Prioridad:} Must \quad \textbf{Estado:} IMPLEMENTED \\[0.2cm]
    \textbf{Módulo:} Backend \quad \textbf{Acceso a datos:} SP \\[0.2cm]
    \textbf{Descripción:} El sistema deberá recuperar los resultados detallados de una simulación consultando múltiples tablas mediante un procedimiento almacenado. \\[0.2cm]
    \textbf{Criterio de aceptación:} Al consultar una simulación existente, el sistema compone sus datos relacionados mediante el procedimiento \texttt{sp\_reporte\_simulacion} y retorna un resultado consistente. \\[0.2cm]
    \textbf{Verificación:} Prueba Unitaria (SimulacionServiceTest.java) \quad \textbf{SQL:} db/procs/sp\_reporte\_simulacion.sql
\end{requisito}

\begin{requisito}{RF-04 - Cálculo agregado de promedio de desempeño de conductor}
    \textbf{Tipo:} FUNC \quad \textbf{Prioridad:} Must \quad \textbf{Estado:} IMPLEMENTED \\[0.2cm]
    \textbf{Módulo:} Backend \quad \textbf{Acceso a datos:} SP \\[0.2cm]
    \textbf{Descripción:} El sistema deberá calcular el promedio global de desempeño de un conductor evaluando todas sus simulaciones completadas. \\[0.2cm]
    \textbf{Criterio de aceptación:} Para un conductor con simulaciones, el promedio calculado por el procedimiento \texttt{sp\_calcular\_promedio\_usuario} coincide con los datos persistidos. \\[0.2cm]
    \textbf{Verificación:} Prueba Unitaria (SimulacionServiceTest.java) \quad \textbf{SQL:} db/procs/sp\_calcular\_promedio\_usuario.sql
\end{requisito}

\begin{requisito}{RF-05 - Generación de reporte consolidado diario de actividad}
    \textbf{Tipo:} FUNC \quad \textbf{Prioridad:} Should \quad \textbf{Estado:} IMPLEMENTED \\[0.2cm]
    \textbf{Módulo:} Backend \quad \textbf{Acceso a datos:} SP \\[0.2cm]
    \textbf{Descripción:} El sistema deberá generar un reporte diario que consolide toda la actividad del sistema para la fecha especificada. \\[0.2cm]
    \textbf{Criterio de aceptación:} El reporte diario incluye la actividad correspondiente a la fecha solicitada. \\[0.2cm]
    \textbf{Verificación:} Prueba Unitaria (SimulacionServiceTest.java) \quad \textbf{SQL:} db/procs/sp\_reporte\_actividad\_diaria.sql
\end{requisito}

\begin{requisito}{RF-06 - Actualización masiva de conductores inactivos}
    \textbf{Tipo:} FUNC \quad \textbf{Prioridad:} Should \quad \textbf{Estado:} IMPLEMENTED \\[0.2cm]
    \textbf{Módulo:} Backend \quad \textbf{Acceso a datos:} SP \\[0.2cm]
    \textbf{Descripción:} El sistema deberá actualizar periódicamente el estado de las cuentas de conductor inactivos mediante una operación masiva y dejar traza de auditoría. \\[0.2cm]
    \textbf{Criterio de aceptación:} La operación masiva cambia únicamente usuarios que cumplen el criterio de inactividad (>6 meses), conserva los demás registros y guarda un registro en la tabla de auditoría. \\[0.2cm]
    \textbf{Verificación:} Prueba de Integración (AuthControllerTest.java) \quad \textbf{SQL:} db/procs/sp\_actualizar\_usuarios\_inactivos.sql
\end{requisito}

\begin{requisito}{RF-07 - Validación cruzada de consistencia de reglas de escenario}
    \textbf{Tipo:} FUNC \quad \textbf{Prioridad:} Must \quad \textbf{Estado:} IMPLEMENTED \\[0.2cm]
    \textbf{Módulo:} Backend \quad \textbf{Acceso a datos:} SP \\[0.2cm]
    \textbf{Descripción:} El sistema deberá validar que las configuraciones de los escenarios no contengan reglas contradictorias antes de su guardado. \\[0.2cm]
    \textbf{Criterio de aceptación:} La validación de escenario detecta configuraciones inconsistentes y acepta configuraciones válidas. \\[0.2cm]
    \textbf{Verificación:} Prueba Unitaria (EscenarioServiceTest.java) \quad \textbf{SQL:} db/procs/sp\_validar\_escenario.sql
\end{requisito}

\begin{requisito}{RF-08 - Generación de código secuencial único para certificado}
    \textbf{Tipo:} FUNC \quad \textbf{Prioridad:} Must \quad \textbf{Estado:} IMPLEMENTED \\[0.2cm]
    \textbf{Módulo:} Backend \quad \textbf{Acceso a datos:} SP \\[0.2cm]
    \textbf{Descripción:} El sistema deberá asignar un código de verificación único y secuencial a cada certificado generado. \\[0.2cm]
    \textbf{Criterio de aceptación:} Cada certificado generado recibe un código secuencial único y reproducible dentro de la transacción usando \texttt{sp\_generar\_codigo\_certificado}. \\[0.2cm]
    \textbf{Verificación:} Prueba Unitaria (SimulacionServiceTest.java) \quad \textbf{SQL:} db/procs/sp\_generar\_codigo\_certificado.sql
\end{requisito}

\begin{requisito}{RF-09 - Evaluación Automática de Comportamiento Vial (IA)}
    \textbf{Tipo:} FUNC \quad \textbf{Prioridad:} Must \quad \textbf{Estado:} TESTED \\[0.2cm]
    \textbf{Módulo:} Motor IA \quad \textbf{Acceso a datos:} CRUD-ORM \\[0.2cm]
    \textbf{Descripción:} El sistema deberá evaluar las decisiones del usuario interactuando con un modelo de IA externo (OpenAI) o un motor de reglas local de respaldo, registrando el identificador del modelo utilizado. \\[0.2cm]
    \textbf{Criterio de aceptación:} El sistema clasifica el comportamiento vial y guarda la decisión con el identificador del modelo utilizado (IA, IA\_LOCAL u OPENAI). \\[0.2cm]
    \textbf{Verificación:} Prueba de Integración (FeedbackServiceTest.java) \quad \textbf{SQL:} No aplica
\end{requisito}

\begin{requisito}{RF-10 - Generación de Retroalimentación Textual (IA)}
    \textbf{Tipo:} FUNC \quad \textbf{Prioridad:} Must \quad \textbf{Estado:} TESTED \\[0.2cm]
    \textbf{Módulo:} Motor IA \quad \textbf{Acceso a datos:} CRUD-ORM \\[0.2cm]
    \textbf{Descripción:} El sistema deberá generar retroalimentación textual y recomendaciones basadas en las acciones del conductor durante la simulación. \\[0.2cm]
    \textbf{Criterio de aceptación:} Todo comportamiento vial evaluado como incorrecto o subóptimo genera una recomendación de mejora registrada en la base de datos. \\[0.2cm]
    \textbf{Verificación:} Prueba de Integración (FeedbackServiceTest.java) \quad \textbf{SQL:} No aplica
\end{requisito}

\begin{requisito}{RF-11 - Reglas de Negocio: Certificación y Puntaje}
    \textbf{Tipo:} FUNC \quad \textbf{Prioridad:} Must \quad \textbf{Estado:} IMPLEMENTED \\[0.2cm]
    \textbf{Módulo:} Backend \quad \textbf{Acceso a datos:} SP \\[0.2cm]
    \textbf{Descripción:} El sistema deberá emitir un certificado si y solo si el puntaje calculado en la simulación es mayor o igual a 80 puntos. \\[0.2cm]
    \textbf{Criterio de aceptación:} Las simulaciones con puntaje >= 80 disparan la generación de certificado, de lo contrario se marcan como fallidas. \\[0.2cm]
    \textbf{Verificación:} Análisis de Inspección (sp\_calcular\_puntaje.sql) \quad \textbf{SQL:} db/procs/sp\_calcular\_puntaje.sql
\end{requisito}

\begin{requisito}{RF-12 - Persistencia de Entidades Secundarias}
    \textbf{Tipo:} FUNC \quad \textbf{Prioridad:} Should \quad \textbf{Estado:} IMPLEMENTED \\[0.2cm]
    \textbf{Módulo:} Database \quad \textbf{Acceso a datos:} CRUD-ORM \\[0.2cm]
    \textbf{Descripción:} El sistema deberá registrar y administrar de manera consistente las Reglas de Tránsito, Infracciones, Eventos Viales, Decisiones y Recompensas generadas durante las prácticas. \\[0.2cm]
    \textbf{Criterio de aceptación:} Las 5 entidades de dominio (ReglaTransito, Infraccion, EventoVial, Decision, Recompensa) persisten su estado relacional sin errores de integridad. \\[0.2cm]
    \textbf{Verificación:} Prueba de Integración (SimulacionIntegracionTest.java) \quad \textbf{SQL:} No aplica
\end{requisito}

\newpage
% ==============================
% 4. REQUISITOS NO FUNCIONALES
% ==============================
\section{Requisitos no funcionales}

\begin{requisito}{RNF-01 - Latencia de lectura en caché caliente p95 $\le$ 200ms}
    \textbf{Tipo:} PERF \quad \textbf{Prioridad:} Must \quad \textbf{Estado:} VERIFIED \\[0.2cm]
    \textbf{Módulo:} API \quad \textbf{Acceso a datos:} No aplica \\[0.2cm]
    \textbf{Descripción:} El sistema deberá responder a consultas cacheadas con un percentil 95 de latencia no mayor a 200 milisegundos bajo una carga de 50 usuarios virtuales concurrentes durante 30 segundos. \\[0.2cm]
    \textbf{Criterio de aceptación:} La corrida k6 reporta p95 de caché caliente <= 200 ms y tasa de error inferior a 1\%. Se conserva la salida cruda NDJSON como evidencia en el repositorio. \\[0.2cm]
    \textbf{Verificación:} Prueba de Rendimiento (scripts/k6/load-test.js) \quad \textbf{SQL:} No aplica
\end{requisito}

\begin{requisito}{RNF-02 - Latencia de lectura en frío p95 $\le$ 500ms}
    \textbf{Tipo:} PERF \quad \textbf{Prioridad:} Must \quad \textbf{Estado:} VERIFIED \\[0.2cm]
    \textbf{Módulo:} API \quad \textbf{Acceso a datos:} No aplica \\[0.2cm]
    \textbf{Descripción:} El sistema deberá responder a consultas sin caché previo con un percentil 95 de latencia no mayor a 500 milisegundos bajo carga de 50 VUs por 30s. \\[0.2cm]
    \textbf{Criterio de aceptación:} La corrida k6 reporta p95 de lectura fría <= 500 ms. Se conserva la salida cruda NDJSON. \\[0.2cm]
    \textbf{Verificación:} Prueba de Rendimiento (scripts/k6/load-test.js) \quad \textbf{SQL:} No aplica
\end{requisito}

\begin{requisito}{RNF-03 - Contraseñas hasheadas con BCrypt y factor de costo $\ge$ 10}
    \textbf{Tipo:} SEC \quad \textbf{Prioridad:} Must \quad \textbf{Estado:} VERIFIED \\[0.2cm]
    \textbf{Módulo:} Backend \quad \textbf{Acceso a datos:} No aplica \\[0.2cm]
    \textbf{Descripción:} El sistema deberá hashear todas las contraseñas utilizando el algoritmo BCrypt con un factor de costo mínimo de 10. \\[0.2cm]
    \textbf{Criterio de aceptación:} Las contraseñas se almacenan con BCrypt y el sistema genera hashes con el prefijo correcto correspondiente al factor de costo configurado. \\[0.2cm]
    \textbf{Verificación:} Prueba de Integración (AuthServiceTest.java) \quad \textbf{SQL:} No aplica
\end{requisito}

\begin{requisito}{RNF-04 - Cookie de token con HttpOnly y SameSite=Strict}
    \textbf{Tipo:} SEC \quad \textbf{Prioridad:} Must \quad \textbf{Estado:} VERIFIED \\[0.2cm]
    \textbf{Módulo:} Backend \quad \textbf{Acceso a datos:} No aplica \\[0.2cm]
    \textbf{Descripción:} El sistema deberá enviar el token de sesión a través de una cookie HTTP configurada con los flags HttpOnly y SameSite=Strict. \\[0.2cm]
    \textbf{Criterio de aceptación:} La cookie de respuesta incluye HttpOnly y SameSite=Strict en cualquier entorno. \\[0.2cm]
    \textbf{Verificación:} Prueba de Integración (AuthControllerTest.java) \quad \textbf{SQL:} No aplica
\end{requisito}

\begin{requisito}{RNF-05 - Ausencia total de SQL dinámico en procedimientos}
    \textbf{Tipo:} SEC \quad \textbf{Prioridad:} Must \quad \textbf{Estado:} VERIFIED \\[0.2cm]
    \textbf{Módulo:} Database \quad \textbf{Acceso a datos:} No aplica \\[0.2cm]
    \textbf{Descripción:} El sistema deberá prevenir la inyección SQL en la base de datos evitando el uso de SQL dinámico inseguro. \\[0.2cm]
    \textbf{Criterio de aceptación:} El análisis automático no detecta concatenación de cadenas dentro de bloques de ejecución dinámica (\texttt{EXECUTE} o \texttt{format}). \\[0.2cm]
    \textbf{Verificación:} Análisis Estático (scripts/audit-sql-dynamic.sh) \quad \textbf{SQL:} No aplica
\end{requisito}

\begin{requisito}{RNF-06 - Escala de Usabilidad del Sistema SUS $\ge$ 68 puntos}
    \textbf{Tipo:} USAB \quad \textbf{Prioridad:} Must \quad \textbf{Estado:} VERIFIED \\[0.2cm]
    \textbf{Módulo:} Frontend \quad \textbf{Acceso a datos:} No aplica \\[0.2cm]
    \textbf{Descripción:} El sistema deberá ofrecer una interfaz altamente usable, alcanzando un puntaje SUS sobre la media estándar de la industria. \\[0.2cm]
    \textbf{Criterio de aceptación:} La evaluación empírica SUS con al menos 15 participantes alcanza una puntuación media calculada igual o superior a 68. \\[0.2cm]
    \textbf{Verificación:} Prueba de Usabilidad (docs/mediciones/sus/sus-analysis.md) \quad \textbf{SQL:} No aplica
\end{requisito}

\begin{requisito}{RNF-07 - Degradación controlada (Fallback) del proveedor de IA}
    \textbf{Tipo:} SEC \quad \textbf{Prioridad:} Should \quad \textbf{Estado:} VERIFIED \\[0.2cm]
    \textbf{Módulo:} Backend \quad \textbf{Acceso a datos:} No aplica \\[0.2cm]
    \textbf{Descripción:} El sistema deberá continuar operando mediante un motor de reglas local si el proveedor externo de IA (OpenAI) supera los 2 segundos de latencia o no está disponible. \\[0.2cm]
    \textbf{Criterio de aceptación:} Ante indisponibilidad del servicio OpenAI, el sistema hace fallback al modelo local sin fallar la simulación. \\[0.2cm]
    \textbf{Verificación:} Prueba de Integración (FeedbackServiceTest.java) \quad \textbf{SQL:} No aplica
\end{requisito}

\begin{requisito}{RNF-08 - Privacidad de Datos y Minimización}
    \textbf{Tipo:} SEC \quad \textbf{Prioridad:} Must \quad \textbf{Estado:} TESTED \\[0.2cm]
    \textbf{Módulo:} Backend \quad \textbf{Acceso a datos:} No aplica \\[0.2cm]
    \textbf{Descripción:} El sistema deberá evitar enviar Información de Identificación Personal (PII) al proveedor externo de Inteligencia Artificial. \\[0.2cm]
    \textbf{Criterio de aceptación:} Solo se envían las decisiones y el contexto del escenario a la API de OpenAI, excluyendo nombres, correos o IDs directos. \\[0.2cm]
    \textbf{Verificación:} Prueba Unitaria (FeedbackIaExternaServiceTest.java) \quad \textbf{SQL:} No aplica
\end{requisito}

\begin{requisito}{RNF-09 - Accesibilidad y Compatibilidad (WCAG 2.2)}
    \textbf{Tipo:} USAB \quad \textbf{Prioridad:} Should \quad \textbf{Estado:} VERIFIED \\[0.2cm]
    \textbf{Módulo:} Frontend \quad \textbf{Acceso a datos:} No aplica \\[0.2cm]
    \textbf{Descripción:} El sistema deberá cumplir con las directrices de accesibilidad web WCAG 2.2 nivel AA. \\[0.2cm]
    \textbf{Criterio de aceptación:} La aplicación puede ser navegada completamente con el teclado y cuenta con un ratio de contraste adecuado validado por Lighthouse. \\[0.2cm]
    \textbf{Verificación:} Análisis Estático Automático (Lighthouse) \quad \textbf{SQL:} No aplica
\end{requisito}

\begin{requisito}{RNF-10 - Capacidad y Volumetría}
    \textbf{Tipo:} PERF \quad \textbf{Prioridad:} Should \quad \textbf{Estado:} TESTED \\[0.2cm]
    \textbf{Módulo:} Arquitectura \quad \textbf{Acceso a datos:} No aplica \\[0.2cm]
    \textbf{Descripción:} El sistema deberá soportar hasta 200 usuarios concurrentes y retener el histórico de prácticas de simulaciones por al menos 5 años. \\[0.2cm]
    \textbf{Criterio de aceptación:} El diseño de la base de datos permite retención a largo plazo y las pruebas de carga demuestran estabilidad con la concurrencia objetivo. \\[0.2cm]
    \textbf{Verificación:} Prueba de Carga (scripts/k6/load-test.js) \quad \textbf{SQL:} No aplica
\end{requisito}


\newpage
% ==============================
% 5. HISTORIAS DE USUARIO
% ==============================
\section{Historias de usuario (Casos de Uso)}

\begin{requisito}{US-01 - Registro autónomo de conductor}
    \textbf{Actor principal:} Conductor\\[0.2cm]
    \textbf{Flujo Principal:} 1. El conductor navega al formulario de registro. 2. Completa los campos obligatorios. 3. El sistema valida los datos y registra al conductor en la base de datos. 4. El sistema notifica el éxito de la operación. \\[0.2cm]
    \textbf{Excepciones:} Si el correo ya existe, el sistema muestra un error y detiene el proceso.\\[0.2cm]
    \textbf{Verificación:} Prueba de Componente E2E (app.component.spec.ts) \quad \textbf{SQL:} No aplica
\end{requisito}

\begin{requisito}{US-02 - Inicio de sesión seguro y obtención de cookie}
    \textbf{Actor principal:} Conductor / Administrador\\[0.2cm]
    \textbf{Flujo Principal:} 1. El usuario ingresa credenciales en el login. 2. El sistema verifica el hash BCrypt. 3. El sistema emite un token JWT como cookie. 4. El usuario es redirigido a su panel de control correspondiente. \\[0.2cm]
    \textbf{Excepciones:} Credenciales incorrectas emiten un error "No Autorizado" (401).\\[0.2cm]
    \textbf{Verificación:} Prueba de Componente E2E (app.component.spec.ts) \quad \textbf{SQL:} No aplica
\end{requisito}

\begin{requisito}{US-03 - Exploración interactiva de escenarios viales}
    \textbf{Actor principal:} Conductor\\[0.2cm]
    \textbf{Flujo Principal:} 1. El conductor solicita la lista de escenarios. 2. El sistema recupera los escenarios habilitados. 3. El sistema pagina y muestra los escenarios con sus configuraciones de clima y dificultad. \\[0.2cm]
    \textbf{Excepciones:} Si no hay escenarios disponibles, se muestra un aviso de "Catálogo vacío".\\[0.2cm]
    \textbf{Verificación:} Prueba de Componente E2E (app.component.spec.ts) \quad \textbf{SQL:} No aplica
\end{requisito}

\begin{requisito}{US-04 - Ejecución de simulación y captura de decisiones}
    \textbf{Actor principal:} Conductor\\[0.2cm]
    \textbf{Flujo Principal:} 1. El conductor inicia un escenario. 2. Toma decisiones ante eventos viales. 3. El motor IA evalúa cada decisión. 4. El sistema calcula y consolida los resultados invocando un procedimiento almacenado. \\[0.2cm]
    \textbf{Excepciones:} Si ocurre una desconexión, la simulación se invalida o se reanuda dependiendo del estado.\\[0.2cm]
    \textbf{Verificación:} Prueba de Componente E2E (app.component.spec.ts) \quad \textbf{SQL:} db/procs/sp\_reporte\_simulacion.sql
\end{requisito}

\begin{requisito}{US-05 - Visualización de métricas y feedback en dashboard}
    \textbf{Actor principal:} Conductor\\[0.2cm]
    \textbf{Flujo Principal:} 1. El conductor accede al Dashboard. 2. El sistema recupera las métricas calculadas. 3. Se presenta el promedio global, simulaciones recientes y sugerencias de mejora. \\[0.2cm]
    \textbf{Excepciones:} Datos insuficientes muestran marcadores por defecto (NA).\\[0.2cm]
    \textbf{Verificación:} Prueba de Componente E2E (app.component.spec.ts) \quad \textbf{SQL:} db/procs/sp\_calcular\_promedio\_usuario.sql
\end{requisito}

\begin{requisito}{US-06 - Descarga de certificado con código secuencial}
    \textbf{Actor principal:} Conductor\\[0.2cm]
    \textbf{Flujo Principal:} 1. El conductor aprueba una simulación (>= 80 pts). 2. El sistema genera un certificado invocando la base de datos. 3. El conductor visualiza el certificado y lo exporta en formato inmutable. \\[0.2cm]
    \textbf{Excepciones:} Si la práctica no fue aprobatoria, el botón de descarga se mantiene oculto.\\[0.2cm]
    \textbf{Verificación:} Prueba de Componente E2E (app.component.spec.ts) \quad \textbf{SQL:} db/procs/sp\_generar\_codigo\_certificado.sql
\end{requisito}

\newpage
% ==============================
% 6. CRITERIOS DE ACEPTACIÓN
% ==============================
\section{Criterios de aceptación de la versión}
\begin{itemize}[label=\textbf{--}]
    \item \textbf{AC-01.} El sistema completo se construye y levanta desde una clonación limpia mediante \texttt{make all}.
    \item \textbf{AC-02.} Las pruebas automatizadas superan la cobertura global de 70\% y el análisis estático no reporta vulnerabilidades graves.
    \item \textbf{AC-03.} Los procedimientos almacenados están versionados, catalogados, parametrizados y trazados.
    \item \textbf{AC-04.} Cinco corridas k6 cumplen los umbrales frío y caliente; Lighthouse y ZAP conservan evidencia reproducible.
    \item \textbf{AC-05.} El despliegue público dispone de healthcheck, runbook, backup y cuenta demo no administrativa.
    \item \textbf{AC-06.} El repositorio contiene documentación, datos, metadatos y etiqueta v1.1.0 coherentes con el commit declarado.
\end{itemize}

\vspace{1cm}
\noindent\textbf{\Large Fin de la especificación}

\end{document}
